\chapter{Proceso de desarrollo, defectos encontrados y decisiones}
\label{chap:desarrollo}

Los capítulos anteriores presentan el protocolo como si hubiera existido desde el principio. No fue
así. El diseño que se ha descrito es el resultado de una secuencia de correcciones, y varias de
ellas nacieron de descubrir que una medición previa era inválida. Este capítulo documenta esa
trayectoria a partir de la bitácora del proyecto.

Hay dos razones para incluirlo en lugar de presentar sólo el resultado final. La primera es de
honestidad: algunos de los defectos que se describen afectaron a cifras que llegaron a documentarse
como resultados, y omitirlos daría una imagen de linealidad que no corresponde a los hechos. La
segunda es que los propios defectos son evidencia. Un sistema que detecta que una de sus puertas
estadísticas penalizaba a los candidatos superiores, y que lo detecta porque ningún retador
resultaba elegible pese a dominar, está mostrando algo sobre su capacidad de rechazar estados
inválidos. Las lecciones son además transferibles: casi ninguna es específica de este dominio.

\section{El estudio como unidad científica}

El diseño original separaba la investigación en dos fases, una exploratoria y otra confirmatoria,
con hipótesis formuladas entre ambas. Se eliminó en julio de 2026 y se sustituyó por un único
\textit{Model Study} automático.

El motivo inmediato fue operativo. Un estudio iniciado el 24 de julio completó un ajuste costoso y
falló después, al consultar un campo de configuración que ya había sido eliminado del código. El
error vivía sólo en memoria, y el proceso desapareció sin dejar reconciliación posible: no había
forma de saber, desde los artefactos, qué se había llegado a calcular. El problema de fondo era que
la arquitectura admitía demasiadas rutas --- escenarios, experimentos, ejecuciones sueltas ---, y
cada ruta era una oportunidad para que un estado parcial quedara sin registrar.

La reconstrucción impuso una regla que el resto del trabajo hereda: existe un único tipo de unidad
científica, se ejecuta entera o falla, y todo lo que produce queda persistido con identidad propia.
Los estados posibles de una ejecución se hicieron explícitos, se añadió un latido para detectar
procesos muertos y una reconciliación por identificador de proceso al arrancar. La consecuencia
práctica para este documento es que cada cifra del Capítulo~\ref{chap:resultados} puede rastrearse
hasta un fichero concreto de un estudio concreto.

Durante la validación de esa reconstrucción aparecieron ocho incidencias que conviene mencionar
porque ilustran cuántas formas hay de que una ejecución «funcione» sin ser válida. La más
instructiva es la segunda: el primer \textit{smoke} técnicamente exitoso produjo puntuaciones
constantes. La causa era que el mínimo de cincuenta observaciones por hoja impedía que los árboles
dividieran un conjunto de sesenta y cinco filas, de modo que el modelo devolvía siempre el mismo
valor. El proceso terminó sin error, escribió sus artefactos y se declaró correcto. No se aceptó
como validación. Es un recordatorio de que la ausencia de excepción no es evidencia de corrección.

\section{Variables declaradas que no se calculaban}

El hallazgo científicamente más relevante de la limpieza de julio fue de otra naturaleza. El
catálogo de \textit{features} declaraba dos bloques de \textit{momentum} multi-horizonte, presentes
en todos los conjuntos de variables reales. Pero esas columnas sólo se calculaban si dos
interruptores concretos estaban activados, y esos dos interruptores \emph{nunca existieron en el
catálogo cerrado}: no eran alcanzables desde ningún estudio real.

La consecuencia es que un estudio previamente documentado se ejecutó sin las columnas de
\textit{momentum} multi-horizonte que su propia configuración declaraba usar. Nada falló, ninguna
prueba se puso en rojo y los resultados eran perfectamente plausibles. El defecto era invisible
porque el sistema declaraba una cosa y computaba otra, y ambas descripciones eran internamente
coherentes.

La lección --- declarar no es computar --- motivó una comprobación que hoy forma parte de la
batería de contratos: que el catálogo declarado y el conjunto de columnas efectivamente calculadas
coincidan. En la misma limpieza se retiró un sesgo de régimen fijado por constantes en el código, no
aprendido de los datos, que contradecía frontalmente el propósito del trabajo.

\section{Cuando la puerta estadística penalizaba al mejor}

La corrección de validez de finales de julio es la más extensa del registro: doce defectos
identificados y medidos, recogidos en la Tabla~\ref{tab:defectos}.

\begin{table}[H]
  \centering
  \footnotesize
  \caption{Defectos de validez detectados y efecto medido de su corrección.}
  \label{tab:defectos}
  % Fuente: docs/bitacora.md, entrada 2026-07-28 «Corrección de validez, alfa neto y limpieza».
% Tabla transcrita del registro de desarrollo, no generada por script: su origen es la
% bitácora del proyecto y no un artefacto del Study.
\begin{tabularx}{\textwidth}{XX}
\toprule
Defecto detectado & Efecto medido de la corrección \\
\midrule
La puerta de no inferioridad penalizaba a los candidatos \textbf{superiores}: cuanto mayor la
diferencia, más ancho el intervalo &
\texttt{feature\_preset} pasa de \texttt{core} (Rank-IC 0,0730) a \texttt{all} (0,0958), con
ventaja pareada $+0{,}0216$. Ningún retador era elegible pese a dominar \\
\addlinespace
\texttt{market\_regime\_feature} se decidió sobre una ventaja de 0,00112, por debajo del ruido &
Pasa de \texttt{True} a \texttt{False} por simplicidad \\
\addlinespace
\texttt{meta\_method} se decidió sobre una ventaja de 0,00033 &
Se mantiene, pero registrado como empate técnico y no como victoria \\
\addlinespace
Emparejamiento por cadena de fecha con rejillas desplazadas por \texttt{execution\_lag\_days} &
Se empareja por periodo mensual; sin bloque completo común el resultado se marca no aplicable en
vez de devolver \texttt{ci\_low = 0,0} \\
\addlinespace
\texttt{evaluation\_key} no incluía el hash del dataset &
La misma clave aparecía con CAGR 0,1468 y 0,1692. Ahora la clave separa datasets \\
\addlinespace
Seis factores de precio se inyectaban en el agente \textit{momentum} fuera de todo condicional &
La ablación \texttt{fundamental} deja de recibir información de precio y mide lo que declara \\
\addlinespace
Dos definiciones incompatibles de \textit{information ratio} bajo el mismo nombre &
Una sola, anualizada \\
\addlinespace
CAGR, \textit{drawdown} y alfa mezclaban la era reservada con la de selección &
Métricas segmentadas: selección, confirmación y curva completa \\
\addlinespace
\texttt{geometric\_excess\_return} era una resta de CAGR &
Cociente de acumulados \\
\addlinespace
Una posición sin precio se marcaba plana &
Convención de exclusión tipo CRSP ($-30\,\%$) y liquidación contra efectivo \\
\addlinespace
\texttt{subsample=0.8} sin \texttt{subsample\_freq}: el \textit{bagging} nunca se activaba &
\texttt{subsample\_freq=1} \\
\addlinespace
Nulo de carteras aleatorias con percentil 95 de CAGR del 107\,\% &
Exige cobertura anual completa, aplica la guarda de datos y paga los mismos costes \\
\bottomrule
\end{tabularx}

  \caption*{Fuente: \artefacto{docs/bitacora.md}, entrada de 2026-07-28.}
\end{table}

El primero de la lista merece detalle porque afecta al corazón del protocolo. La regla de selección
admite un candidato si demuestra no ser inferior al incumbente, y esa no inferioridad se evaluaba
con un intervalo de confianza sobre la diferencia pareada. El defecto era que el intervalo se
ensanchaba con la magnitud de la diferencia: cuanto \emph{mejor} era un retador, más ancho su
intervalo y más difícil superar el umbral. La puerta penalizaba exactamente a los candidatos que
debía premiar.

El síntoma que lo delató no fue un error, sino una regularidad sospechosa: ningún retador resultaba
elegible, en ninguna variable, pese a que algunos dominaban con claridad. Al corregirlo,
\texttt{feature\_preset} pasó de \texttt{core} (Rank-IC 0,0730) a \texttt{all} (0,0958), con una
ventaja pareada de $+0{,}0216$. Es decir, el sistema había estado descartando una mejora
sustancial y real por un defecto en su criterio de admisión.

Dos defectos más de esa lista tienen implicaciones directas para la lectura del
Capítulo~\ref{chap:resultados}. La clave de evaluación no incluía el identificador del conjunto de
datos, de modo que la misma clave llegó a aparecer asociada a CAGR 0,1468 y 0,1692: dos resultados
distintos indexados como si fueran el mismo. Y el exceso geométrico se calculaba como una resta de
CAGR en lugar de un cociente de acumulados, lo que introduce un sesgo que crece con el horizonte.
Ambas correcciones cambian cifras publicadas, y por eso los estudios anteriores a esta versión del
catálogo quedaron derogados como evidencia.

De la misma revisión salió una observación que motivó introducir un conjunto de semillas: el Rank-IC
variaba apenas $\pm 0{,}001$ entre inicializaciones, pero el exceso sobre el índice \emph{cambiaba
de signo} (una semilla daba $-0{,}51$ puntos y otra $+3{,}11$). La conclusión predictiva era
estable; la económica, no. Esa asimetría es la que la Figura~\ref{fig:semillas} documenta hoy con
tres semillas.

\section{Toda venta necesita un destino mejor}

El rediseño de la cartera partió de tres defectos que, vistos por separado, parecían independientes,
y que compartían una única raíz: las ventas se decidían sin mirar a dónde iba el dinero.

El primero era el más grave en términos contables. Bajo la política de efectivo oportunista, si
ningún candidato superaba el umbral, la lógica vendía la cartera entera y devolvía una cartera
objetivo vacía. El \textit{backtest} interpretaba «vacía» como «mantener» y resucitaba las
posiciones --- pero ya había cargado los costes de unas ventas que nunca se ejecutaron. La cartera
pagaba por replegarse sin replegarse.

El segundo producía el efecto contrario y contribuyó de forma decisiva al 877\,\% de rotación anual
que se observaba entonces. Bajo la política de permanecer siempre invertido, las doce posiciones se
vendían por umbral y el relleno obligatorio \emph{recompraba exactamente las mismas en el mismo
snapshot}: una ida y vuelta completa, con sus costes, para acabar donde se estaba.

El tercero era la ausencia de un límite de concentración: con un único candidato admisible y un tope
de efectivo del 20\,\%, el 80\,\% de la cartera acababa en un solo valor.

La regla que los resuelve a los tres es la que gobierna hoy la cartera: una venta sólo se emite si
el destino del dinero es mejor que la posición, después de costes. El destino puede ser otro valor
--- y entonces debe superar el umbral de rotación --- o el efectivo, que exige un tope declarado y
respeta un suelo de diversificación. Se añadió además histéresis entre el umbral de entrada y el de
salida, para que una posición no oscile dentro y fuera por variaciones mínimas.

\section{Umbrales anuales y comparabilidad}

Un cambio aparentemente menor tiene consecuencias sobre todo lo demás. Los umbrales de la cartera se
expresaban en puntos básicos referidos al horizonte, no al año. Con esa convención, 250 puntos
básicos en un horizonte de tres meses equivalen aproximadamente a un 10\,\% anual, mientras que en
uno de doce meses equivalen a un 2,5\,\%: el mismo número significaba exigencias muy distintas según
otra variable del catálogo.

La corrección fue expresar todos los umbrales en puntos básicos anuales y convertirlos al horizonte
mediante composición geométrica. El prorrateo lineal se rechazó de forma deliberada: es el mismo
atajo aritmético que ya había sido corregido en el CAGR y en el \textit{information ratio}, y
reintroducirlo por comodidad habría sido inconsistente.

\section{El colapso de la calibración isotónica}

En julio de 2026 se observó un comportamiento que no encajaba: posiciones con ochenta y muchos meses
de antigüedad y percentil actual muy bajo que nunca se vendían. El diagnóstico encadenó dos
problemas distintos.

El primero era que el meta-agente no olvidaba el régimen antiguo. Su Rank-IC acumuló siete
trimestres negativos consecutivos, una racha más larga que la de la crisis sanitaria, y los cinco
agentes se invirtieron a la vez --- un patrón compatible con una rotación de factores, no con un
error de signo. El ajuste usaba una ventana rectangular: una cohorte de cuatro años atrás pesaba lo
mismo que la más reciente.

El segundo era más sutil y bloqueaba la cartera por completo. La calibración de percentil a alfa
esperada usaba una regresión isotónica creciente. Cuando la relación real entre percentil y retorno
futuro es \emph{decreciente}, la única curva creciente que minimiza el error es una constante. El
resultado fue que la alfa esperada era idéntica para los quinientos cuatro valores del universo. Con
una alfa plana no existe ventaja calculable entre dos posiciones, y el bucle de rotación se
interrumpía sin emitir ninguna orden. No era un fallo de implementación: era el comportamiento
correcto de un estimador aplicado fuera de su supuesto.

El diagnóstico de ese episodio contiene una lección estadística que merece registrarse por sí misma.
El primer análisis por deciles, hecho con la \emph{media} del retorno excedente, sugería que la
ordenación estaba invertida en todo el histórico. Era un artefacto del estimador: la distribución
tiene una cola derecha extrema, y unos pocos valores multiplicadores situados en el decil bajo
bastaban para que su media (${+}3552$ puntos básicos) contradijera a su mediana (${-}146$). Con la
mediana, o con una media winsorizada, la curva histórica \emph{sí} es creciente, de forma coherente
con el Rank-IC positivo. La inversión era real, pero sólo en el régimen reciente. Un diagnóstico
apresurado habría llevado a reescribir el sistema entero para corregir un problema que no existía.

La isotónica se sustituyó por una recta de percentil a retorno anualizado real con una cascada de
ventanas, que acepta la primera ventana con pendiente creciente. La estimación se hace sobre veinte
ventiles, pero la evaluación se realiza en el rango continuo: produce quinientos cuatro valores de
alfa distintos donde la isotónica producía uno.

\section{Dos decisiones tomadas contra la recomendación técnica}

El registro documenta dos decisiones que se adoptaron en contra de la recomendación técnica. Se
recogen aquí con ambos lados del argumento, porque omitirlas presentaría el protocolo como más
limpio de lo que es.

La primera fue sustituir la calibración isotónica sin conservarla como opción del catálogo. A favor:
el colapso es demostrable y deja la cartera inoperativa, de modo que mantener la opción habría
significado conservar una rama que se sabe rota. En contra: cambia el ganador sin pasar por la
selección secuencial por Rank-IC, que es precisamente la regla 4 del protocolo, y renuncia a poder
comparar ambas formulaciones con evidencia. El autor asumió el riesgo y pidió eliminar la
alternativa.

La segunda es más incómoda. Cuando ninguna de las tres ventanas de la cascada produce una pendiente
creciente, el sistema aplica una salvaguarda: una recta fija entre $-10\,\%$ y $+10\,\%$. En contra
de esa elección se argumentó que es un supuesto \emph{a priori}, del mismo tipo que el sesgo de
régimen fijado por constantes que este proyecto ya había eliminado, y que se activa precisamente
cuando las tres ventanas coinciden en que la ordenación no discrimina a favor --- es decir, cuando
la evidencia disponible dice que no hay señal, el sistema inventa una. La recomendación técnica fue
usar una curva plana, equivalente a repartir por igual. Se prefirió la recta fija.

La mitigación consiste en registrar, para cada fila, qué ventana produjo su alfa. En el estudio de
referencia la salvaguarda se activa en 4.384 de 20.545 filas, algo más de una de cada cinco. Es una
proporción suficientemente alta como para que el lector deba tenerla presente al interpretar
cualquier cifra de cartera, y por eso se recoge también en el Capítulo~\ref{chap:limitaciones}.

Durante esa misma revisión se descubrió que la ruta de calibración de alfa \emph{no tenía ninguna
cobertura de pruebas}. Un renombrado había dejado una llamada rota en esa ruta que la batería
existente no detectaba. Se añadió un contrato específico que cubre justamente lo que puede fallar en
silencio: que cada rango produzca una alfa distinta y monótona, que una relación decreciente caiga a
la salvaguarda, y que la ausencia de cohortes cerradas produzca un valor no disponible en lugar de
un cero.

\section{Una posición sin puntuación detenía toda la rotación}

Un episodio concreto ilustra por qué las casuísticas de datos incompletos merecen tratamiento
explícito. En una ejecución, una posición estaba en el percentil 4,96 con un retorno esperado
de $-9{,}01\,\%$ anual, mientras otro valor fuera de la cartera estaba en el percentil 100 con
$+10{,}00\,\%$: una ventaja de 1.901 puntos básicos, muy por encima de cualquier umbral de rotación.
El intercambio no se produjo.

La causa era que una tercera posición, de una empresa que había salido del índice años antes, no
tenía fila puntuable en esa fecha. El bucle la seleccionaba como peor posición, no podía calcular su
ventaja y se detenía --- bloqueando de paso todas las rotaciones posteriores, incluida la que sí era
claramente rentable. La corrección fue tratar la ausencia de puntuación como un mandato de venta,
no como una decisión económica: una posición sin puntuación actual se vende, ignorando el mínimo de
tenencia, y no se recompra en ese mismo snapshot.

\section{El libro de versiones del catálogo}

Cada uno de estos cambios rompe la comparabilidad con los estudios anteriores, y el proyecto lleva
un registro explícito de esas rupturas mediante una versión de catálogo. La
Tabla~\ref{tab:versiones} lo resume junto con el crecimiento de la batería de contratos.

\begin{table}[H]
  \centering
  \footnotesize
  \caption{Versiones del catálogo cerrado y crecimiento de la batería de pruebas.}
  \label{tab:versiones}
  % Fuente: docs/bitacora.md (entradas 2026-07-25 a 2026-08-04) y module/studies/catalog.py.
% Transcrita del registro de desarrollo, no generada por script.
\begin{tabularx}{\textwidth}{clXr}
\toprule
Versión & Fecha & Qué cambia y por qué rompe la comparabilidad & Tests \\
\midrule
2 & 2026-07-25 &
Reconstrucción a Model Study único: se elimina el protocolo exploratorio--confirmatorio y se retiran
las variantes de meta con sesgo no aprendido (\texttt{regime}, \texttt{rank\_ic}) &
15 \\
\addlinespace
3 & 2026-07-28 &
Corrección de doce defectos de validez; los umbrales de cartera pasan de percentiles a puntos
básicos y \texttt{sizing\_mode} de \texttt{score\_linear} a \texttt{alpha\_proportional} &
--- \\
\addlinespace
4 & 2026-07-28 &
Rediseño de la gestión de cartera bajo el principio de que toda venta necesita un destino mejor;
\texttt{max\_cash\_weight} entra con rejilla $(0\,\%,\,10\,\%,\,25\,\%)$ &
45 \\
\addlinespace
5 & 2026-07-28 &
Umbrales expresados en puntos básicos anuales con conversión geométrica; entran
\texttt{minimum\_holding\_period} y \texttt{price\_only\_sell\_only} &
52 \\
\addlinespace
6 & 2026-08-04 &
\texttt{cash\_policy} se elimina y \texttt{max\_cash\_weight} gobierna sola el efectivo; se corrige
el reparto que dejaba efectivo no declarado y entra \texttt{coverage\_percentile\_floor} &
75 \\
\bottomrule
\end{tabularx}

  \caption*{Fuente: \artefacto{docs/bitacora.md} y \artefacto{module/studies/catalog.py}.}
\end{table}

El número de contratos pasó de 15 a 75 a lo largo del desarrollo, y casi todos los añadidos
proceden de un defecto concreto: cada corrección de esta lista dejó tras de sí una prueba que
impide su reaparición. Esa es, probablemente, la parte más reutilizable del trabajo.

\section{Advertencia de reproducibilidad}
\label{sec:reproducibilidad}

Hay una consecuencia de este historial que debe declararse sin rodeos. El estudio del que procede
toda la evidencia de este trabajo, \studyid, se ejecutó bajo la versión 5 del catálogo. El código
actual está en la versión 6, tras la corrección del reparto de efectivo del 4 de agosto de 2026.

Los artefactos siguen siendo válidos y trazables: conservan sus identificadores, sus hashes de
conjunto de datos y su instantánea de catálogo, y cualquier cifra de este documento puede
verificarse contra ellos. Pero \textbf{no son reproducibles bit a bit con el código actual}, y no
son comparables con estudios ejecutados bajo la versión 6.

La alternativa habría sido reejecutar el estudio completo. Se descartó por una razón metodológica y
no de conveniencia: la reejecución cambiaría el ganador, y escribir los capítulos de resultados
sobre un ganador y después sustituirlo por otro cuando el texto ya está redactado se parece
demasiado a elegir la configuración cuyos resultados encajan mejor con lo escrito. Se optó por
congelar el estudio, declarar la discrepancia y mantener todas las cifras referidas a la versión con
la que realmente se calcularon.
